\chapter{共享对象的事件驱动回收}
\label{ch:reclamation}

\section{索引移除与读者访问之间的时间差}
共享对象移出索引后，已经获得引用的读者仍可能继续访问其存储。若索引移除立即触发释放，后续读者访问便可能指向已经回收的空间；若始终保留退役对象，又会延长存储占用。本章在教学模型中研究何时可以安全释放退役对象，并比较事件检查与周期扫描的决策时点。引用计数、退役和回收沿用第\ref{ch:observation}章的定义。

\section{状态与安全条件}
对象依次经历在索引中、已退役和已回收三种状态。对象退役后不再接受新引用，既有引用则随读者退出而减少。记退役状态为 $retired_i$，回收的必要条件为
\begin{equation}
retired_i \land (r_i=0).
\label{eq:reclaim-condition}
\end{equation}
退役条件阻止新读者进入，零引用条件排除尚在访问的读者。二者共同成立才允许释放空间。该模型假设获取引用、退役、引用更新以及检查后释放均受同一把互斥锁保护；每个读者结束访问后恰好释放一次引用。因此，状态检查与实际释放之间不会插入新的引用获取。

这两项条件分别排除了不同的不安全状态。在索引中的对象即使暂时没有读者，也可能被随后到来的读者获取；已退役对象即使不能再被新读者发现，也可能仍被既有读者持有。由模型假设可知，满足式\ref{eq:reclaim-condition}时，既无当前访问者，也无合法的新引用获取路径。这给出了该教学状态机内释放操作的安全性依据。

\section{贯穿示例与事件处理}
考虑初始时仍在索引中的对象 $x$，有两个读者持有它，即 $r_x=2$。退役操作首先移除索引项，将对象状态改为已退役，并调用回收检查。此时引用计数仍为2，存储继续保留。第一个读者退出后，引用计数降为1，检查仍不能释放对象。第二个读者退出后，引用计数降为0；由于对象已经退役，此次检查释放其存储并将状态改为已回收。该示例描述事件次序，没有设定各事件间的时间间隔。

伪代码中的退役处理与读者退出处理均调用同一个检查过程。退役时调用检查，可以覆盖对象移出索引时已经没有读者的情形；读者退出时调用检查，可以覆盖最后一个读者在退役之后才退出的情形。统一检查过程同时测试状态和计数，避免两个事件处理路径采用不同条件。状态变为已回收后，不再执行重复释放。

\section{决策时点与实现代价}
周期扫描策略也采用式\ref{eq:reclaim-condition}，但仅在扫描时发现满足条件的对象。事件检查把判断放在可能使条件成立的退役和引用释放事件之后，因而在该状态机中无需等待下一次扫描。两种策略共享安全条件，差别在于何时检查条件。

在句柄定位按常数成本计算的模型下，单次事件检查只读取对象状态和引用计数，计算成本为 $O(1)$；一次周期扫描需检查退役集合中的 $q$ 个对象，成本为 $O(q)$。这两个复杂度对应不同触发频率，不能仅凭单次成本判断总成本。事件检查把额外操作放入退役和读者退出的执行路径，还需保存状态与计数，并在本模型中取得互斥锁。实际成本取决于事件数量和竞争情况。

本章沿用前章的对象与计数定义。回收器根据对象的退役状态和读者引用计数作出判断，不调用前章的观测工具。两章分别研究观测描述与回收决策，联系在于共享同一对象模型。

\section{合成结果回答的问题}
表\ref{tab:synthetic-space}列出用于教学分析的合成值，比较相同轨迹及对象输入下两种策略的峰值保留空间。统计范围包括对象存储及各自的策略元数据。
\begin{table}[ht]
\centering
\caption{教学合成值：对象存储与策略元数据的峰值总和（KiB）。}
\label{tab:synthetic-space}
\begin{tabular}{lrr}
轨迹 & 周期扫描 & 事件检查 \\
W & 100 & 72 \\
V & 100 & 110
\end{tabular}
\end{table}

在轨迹W中，事件检查的峰值保留空间为72 KiB，相比周期扫描的100 KiB降低28\%；在轨迹V中，事件检查为110 KiB，相比100 KiB增加10\%。这组对照回答的是不同轨迹下总体空间量如何变化，呈现了收益随输入而变化的情形。表中只有总空间，不能分离回收时点与元数据各自的影响，也不能据此判断延迟或吞吐率。

\section{本章小结}
本章围绕退役与读者退出的不同步，说明共享对象回收所需的两个条件。事件驱动检查在退役和引用释放后判断条件，通过明确的状态转换解释释放时点；互斥保护和完整引用计数构成该教学模型的安全前提。合成对照展示了不同轨迹下峰值空间可能下降或增加，结果解释保留了指标与条件。
